\section{Equivalence between direct sum theorems and XOR lemmas and the proof of~\Cref{thm:main-xor}}

In this section, we prove the following claim which shows that a strong direct sum theorem implies a strong XOR lemma. We then derive~\Cref{thm:main-xor} as a consequence of this equivalence and our strong direct sum theorem for query complexity (\Cref{thm:main-intro}).

% \lnote{Add prose. Right now~\Cref{claim:equivalence between direct product and XOR} comes out of nowhere.}
\begin{claim}[Equivalence between direct sum theorems and XOR lemmas]
\label{claim:equivalence between direct product and XOR}
For every $f:\bn\to\bits$, distribution $\mu$ over $\bn$, integer $k\in\N$, multiplicative factor $M \in \R$, and $\eps\in (0,1)$, if the following direct sum theorem holds:
$$
\overline{\mathrm{Depth}}^{\mu^k}(f^{\otimes k},\eps)\ge M\cdot\overline{\mathrm{Depth}}^{\mu}(f,\delta).
$$
then, the following XOR lemma holds:
$$
\overline{\mathrm{Depth}}^{\mu^k}(f^{\oplus k},\tfrac{\eps}{2})\ge M\cdot\overline{\mathrm{Depth}}^{\mu}(f,\delta).
$$
\end{claim}

In order to prove \Cref{claim:equivalence between direct product and XOR} and \Cref{thm:main-xor}, we establish a lemma which allows us to convert any decision accurately computing $f^{\oplus k}$ into a decision tree accurately computing $f^{\otimes k}$. The following definition captures this conversion.

\begin{definition}[The product tree]
\label{def:product tree}
    Given a decision tree $T:(\bn)^k \to \bits$, the $k$-wise product tree $\Tilde{T}:(\bn)^k \to \bits^k$ is defined as follows. For the internal nodes, $\Tilde{T}$ has exactly the same structure as $T$. For a leaf $\ell$ in $T$, the leaf vector $(\ell_1,\ldots,\ell_k)\in\bits^k$ in $\Tilde{T}$ is defined by
    $$
    \ell_i\coloneqq\sign\left(\Ex_{\bX\sim\mu^{k}}[f(\bX^{(i)})\mid \bX\text{ reaches }\ell]\right)
    $$
    for all $i\in[k]$.
\end{definition}

Intuitively, $\Tilde{T}$ computes $T$'s best guess for $f(X^{(i)})$ for each $i\in [k]$ on a given input $(X^{(1)},\ldots,X^{(k)})$. If $T$ is really good at computing $f^{\oplus k}$ then at every leaf it should have queried enough variables to pin down $f$'s value on each of the input blocks. The main lemma formalizes this intuition.

\begin{lemma}
\label{lem:dpt-to-ds}
    For any $f:\bits^n \to \bits$, distribution $\mu$ over $\bits^n$, and tree $T:(\bn)^k \to \bits$, the $k$-wise product tree $\Tilde{T}:(\bn)^k \to \bits^k$ satisfies
    \begin{equation*}
        \Prx_{\bX \sim \mu^{k}}[\tilde{T}(\bX) = f^{\otimes k}(\bX)] \geq \Ex_{\bX \sim \mu^{k}}[T(\bX)f^{\oplus k}(\bX)].
    \end{equation*}
\end{lemma}



\subsection{Proofs of~\Cref{claim:equivalence between direct product and XOR} and~\Cref{thm:main-xor} assuming~\Cref{lem:dpt-to-ds}}
The following corollary of \Cref{lem:dpt-to-ds} implies~\Cref{claim:equivalence between direct product and XOR} and~\Cref{thm:main-xor}.
\begin{corollary}[Main corollary of \Cref{lem:dpt-to-ds}]
    \label{cor:xor-to-dp-for-r}
    For all $f:\bits^n\to\bits$, $k\ge 1$, distributions $\mu$ over $\bits^n$, and $\eps>0$, $\overline{\mathrm{Depth}}^{\mu^{k}}(f^{\oplus k},\frac{\eps}{2})\ge \overline{\mathrm{Depth}}^{\mu^{k}}(f^{\otimes k},\eps)$.
\end{corollary}
\begin{proof}
    Let $\bA$ be a randomized query algorithm for $f^{\oplus k}$ with error $\eps/2$ and expected cost $q=\overline{\mathrm{Depth}}^{\mu^{k}}(f^{\oplus k},\eps/2)$. Let $\mathcal{T}$ denote the distribution over decision trees determined by $\bA$. Consider the algorithm $\Tilde{\bA}$ which computes $f^{\otimes k}(X)$ by sampling $\bT\sim\mathcal{T}$ and returning $\Tilde{\bT}(X)$ where $\Tilde{\bT}$ is the decision tree from \Cref{lem:dpt-to-ds}. Then, the success of $\Tilde{\bA}$ is
    \begin{align*}
        \Ex_{\bT\sim \mathcal{T}}\left[\Prx_{\bX\sim \mu^{k}}[\Tilde{\bT}(\bX)=f^{\otimes k}(\bX)]\right]&\ge \Ex_{\bT\sim \mathcal{T}}\left[\Ex_{\bX\sim\mu^{k}}[\bT(\bX)f^{\oplus k}(\bX)]\right]\tag{\Cref{lem:dpt-to-ds}}\\
        &\ge 1-\eps
    \end{align*}
    where the last step uses the fact that advantage is $1-2\cdot\mathrm{error}$. Since the structure of each $\Tilde{T}$ is the same as $T$, the expected cost of $\Tilde{\bA}$ is $q$ which completes the proof.
\end{proof}

\begin{proof}[Proofs of~\Cref{claim:equivalence between direct product and XOR} and~\Cref{thm:main-xor}]
    By \Cref{cor:xor-to-dp-for-r}, $$\overline{\mathrm{Depth}}^{\mu^{k}}(f^{\oplus k},\tfrac{\eps}{2})\ge \overline{\mathrm{Depth}}^{\mu^{k}}(f^{\otimes k},\eps)\ge M\cdot\overline{\mathrm{Depth}}^{\mu}(f,\delta).$$
    \Cref{thm:main-xor} follows immediately by applying \Cref{claim:equivalence between direct product and XOR} to \Cref{thm:main-formal}.
\end{proof}

\begin{remark}[On the necessity of the $1/2$ loss in $\eps$ in \Cref{cor:xor-to-dp-for-r}]
    \label{rem:loss}
    One may wonder whether the $1/2$ loss in $\eps$ parameter in \Cref{cor:xor-to-dp-for-r} is necessary. For example, can one show $\overline{\mathrm{Depth}}^{\mu^{k}}(f^{\oplus k},0.51\eps)\ge \overline{\mathrm{Depth}}^{\mu^{k}}(f^{\otimes k},\eps)$? The issue is that $\overline{\mathrm{Depth}}^{\mu^{k}}(f^{\oplus k},0.5)=0$ for all functions $f:\bn\to\bits$ because the bias of $f^{\oplus k}$ is at least $0.5$. So such a statement cannot hold for all $f$ in all parameter regimes. Concretely, one can show that if $f$ is the parity of $n$ bits and $\mu$ is uniform over $\bn$, then $\overline{\mathrm{Depth}}^{\mu^{k}}(f^{\otimes k},\eps)\ge \Omega(kn)$ for all constant $\eps<1$. Any path in a decision tree for $f^{\otimes k}$ which queries at most $\lambda kn$ bits for some constant $\lambda<1$ has success probability $2^{-\Omega(k)}$. So to achieve any constant accuracy requires $\Omega(kn)$ expected depth. On the other hand, $\overline{\mathrm{Depth}}^{\mu^{k}}(f^{\oplus k},0.5)=1\ll \overline{\mathrm{Depth}}^{\mu^{k}}(f^{\otimes k},\eps)$ which is achieved by the decision tree that outputs a single constant value. Therefore, the $\eps/2$ in \Cref{cor:xor-to-dp-for-r} is necessary for such a statement to hold in full generality.
\end{remark}



\subsection{Proof of~\Cref{lem:dpt-to-ds}}

% In this section, we prove the following strong XOR lemma for distributional query complexity.

% \begin{theorem}[Strong XOR lemma for distributional query complexity]
% \label{thm:xor-lemma-technical}
%     For every function $f:\bits^n\to\bits$ and distribution $\mu$ over $\bits^n$, $k\in\N$, and $\gamma,\delta\in(0,1)$, we have that
%     $$
%     \overline{\mathrm{Depth}}^{\mu^{k}}\left(f^{\oplus k},\tfrac{1}{2}(1-(1-\Omega(\delta))^k-\gamma)\right)\ge {\Omega}\left(\frac{\gamma^2 k}{\log(1/\delta)}\right)\cdot\overline{\mathrm{Depth}}^\mu(f,\delta).
%     $$
% \end{theorem}

% Before proving the theorem, we establish a lemma which allows us to convert any decision accurately computing $f^{\oplus k}$ into a decision tree accurately computing $f^{\otimes k}$. The following definition captures this conversion.

% \begin{definition}[The product tree]
% \label{def:product tree}
%     Given a decision tree $T:(\bn)^k \to \bits$, the $k$-wise product tree $\Tilde{T}:(\bn)^k \to \bits^k$ is defined as follows. For the internal nodes, $\Tilde{T}$ has exactly the same structure as $T$. For a leaf $\ell$ in $T$, the leaf vector $(\ell_1,\ldots,\ell_k)\in\bits^k$ in $\Tilde{T}$ is defined by
%     $$
%     \ell_i\coloneqq\sign\left(\Ex_{\bX\sim\mu^{k}}[f(\bX^{(i)})\mid \bX\text{ reaches }\ell]\right)
%     $$
%     for all $i\in[k]$.
% \end{definition}

% Intuitively, $\Tilde{T}$ computes $T$'s best guess for $f(X^{(i)})$ for each $i\in [k]$ on a given input $(X^{(1)},\ldots,X^{(k)})$. If $T$ is really good at computing $f^{\oplus k}$ then at every leaf it should have queried enough variables to pin down $f$'s value on each of the input blocks. The next lemma formalizes this intuition.

% \begin{lemma}
% \label{lem:dpt-to-ds}
%     For any $f:\bits^n \to \bits$, distribution $\mu$ over $\bits^n$, and tree $T:(\bn)^k \to \bits$, the $k$-wise product tree $\Tilde{T}:(\bn)^k \to \bits^k$ satisfies
%     \begin{equation*}
%         \Prx_{\bX \sim \mu^{k}}[\tilde{T}(\bX) = f^{\otimes k}(\bX)] \geq \Ex_{\bX \sim \mu^{k}}[T(\bX)f^{\oplus k}(\bX)].
%     \end{equation*}
% \end{lemma}

Each $\ell_i$ for $i\in [k]$ satisfies
\begin{align*}
    \Ex_{\bX\sim\mu^{k}}[\ell_i\cdot f(\bX^{(i)})\mid\bX\text{ reaches }\ell]&=\left|\Ex_{\bX\sim\mu^{k}}[f(\bX^{(i)})\mid\bX\text{ reaches }\ell]\right|\\
    &\ge \Ex_{\bX\sim\mu^{k}}[\ell\cdot f(\bX^{(i)})\mid\bX\text{ reaches }\ell].
\end{align*}
In particular, for all {leaves} $\ell$ of $T$,
\begin{equation}
\label{eq:corr}
    \Ex_{\bX\sim\mu^{k}}[\Tilde{T}(\bX)_i\cdot f(\bX^{(i)})\mid\bX\text{ reaches }\ell]\ge \Ex_{\bX\sim\mu^{k}}[T(\bX)\cdot f(\bX^{(i)})\mid\bX\text{ reaches }\ell].
\end{equation}

% $$
% \Ex_{\bX\sim\mu^{k}}[\ell_i\cdot f(\bX^{(i)})\mid\bX\text{ reaches }\ell]=\left|\Ex_{\bX\sim\mu^{k}}[f(\bX^{(i)})\mid\bX\text{ reaches }\ell]\right|\ge \Ex_{\bX\sim\mu^{k}}[\ell\cdot f(\bX^{(i)})\mid\bX\text{ reaches }\ell]
% $$
Therefore:
    \begin{align*}
    \Prx_{\bX \sim \mu^{k}}&[\tilde{T}(\bX)= f^{\otimes k}(\bX)]\\
        &= \Ex_{\bell\sim\mu^{k}(T)}\bigg[\Prx_{\bX\sim\mu^{k}}[\tilde{T}(\bX)=f^{\otimes k}(\bX)\mid \bX\text{ reaches }\bell]\bigg]\\
        &=\Ex_{\bell\sim\mu^{k}(T)}\bigg[\prod_{i\in [k]}\Prx_{\bX\sim\mu^{k}}[\tilde{T}(\bX)_i=f(\bX^{(i)})\mid \bX\text{ reaches }\bell]\bigg]\tag{\Cref{claim:leaf-product}}\\
        % &=\Ex_{\bell\sim\mu^{k}(T)}\bigg[\prod_{i\in [k]}\Ex_{\bX\sim\mu^{k}}\bigg[\frac{1+\Tilde{T}(\bX)^{(i)}\cdot f(\bX^{(i)})}{2}\mid \bX\text{ reaches }\bell\bigg]\bigg]
        &\ge \Ex_{\bell\sim\mu^{k}(T)}\bigg[\prod_{i\in [k]}\Ex_{\bX\sim\mu^{k}}[\Tilde{T}(\bX)_if(\bX^{(i)})\mid \bX\text{ reaches }\bell]\bigg]\\
        &\ge \Ex_{\bell\sim\mu^{k}(T)}\bigg[\prod_{i\in [k]}\Ex_{\bX\sim\mu^{k}}[T(\bX)f(\bX^{(i)})\mid \bX\text{ reaches }\bell]\bigg]\tag{\Cref{eq:corr}}\\
        &=\Ex_{\bell\sim\mu^{k}(T)}\bigg[\Ex_{\bX\sim\mu^{k}}\bigg[T(\bX)\prod_{i\in [k]}f(\bX^{(i)})\mid \bX\text{ reaches }\bell\bigg]\bigg]\tag{\Cref{claim:leaf-product}}\\
        % &=\Ex_{\bell\sim\mu^{k}(T)}\bigg[T(\bell)\Ex_{\bX\sim\mu^{k}}\bigg[\prod_{i\in [k]}f(\bX^{(i)})\mid \bX\text{ reaches }\bell\bigg]\bigg]\\
        &=\Ex_{\bell\sim\mu^{k}(T)}\bigg[\Ex_{\bX\sim\mu^{k}}[T(\bX)f^{\otimes k}(\bX)\mid \bX\text{ reaches }\bell]\bigg]\tag{$\prod_{i\in [k]}f(\bX^{(i)})=f^{\oplus k}(\bX)$}\\
        &=\Ex_{\bX \sim \mu^{k}}[T(\bX)f^{\oplus k}(\bX)]
    \end{align*}
    which completes the proof.\hfill\qed

% \begin{proof}\lnote{Minor, but both these proofs are so simple, basically trivial. Can we just replace them with prose?}[Proof of \Cref{claim:equivalence between direct product and XOR}]
%     We have that
%     \begin{align*}
%         \overline{\mathrm{Depth}}^{\mu^{k}}(f^{\oplus k},\tfrac{\eps}{2})&\ge \overline{\mathrm{Depth}}^{\mu^{k}}(f^{\otimes k},\eps)\tag{\Cref{cor:xor-to-dp-for-r}}\\
%         &\ge M\cdot\overline{\mathrm{Depth}}^{\mu}(f,\delta).
%     \end{align*}
% \end{proof}


% With this corollary in hand, we prove the main theorem of this section.

% \begin{proof}[Proof of \Cref{thm:xor-lemma-technical}]
%     We have that
%     \begin{align*}
%         \overline{\mathrm{Depth}}^{\mu^{k}}\left(f^{\oplus k},\frac{1}{2}(1-(1-\Omega(\delta))^k-\gamma)\right)&\ge \overline{\mathrm{Depth}}^{\mu^{k}}\left(f^{\otimes k},1-(1-\Omega(\delta))^k-\gamma\right)\tag{\Cref{cor:xor-to-dp-for-r}}\\
%         &\ge \Omega\left( \frac{\gamma^2 k}{\log(1/\delta)}\right) \cdot \overline{\mathrm{Depth}}^{\mu}(f,\delta)\tag{\Cref{thm:main-formal}}
%     \end{align*}
%     which completes the proof.
% \end{proof}


